% !TeX root = ../papers/nc-p3.tex
\section{種数5の水準の平面四次商と有限非分岐境界}
\label{sec:rank-fifteen-pencil-unramified}
\begin{proposition}\label{prop:rank-fifteen-pencil-plane}
$s=c+c^{-1}$、$s^2\notin\{0,1,4\}$ とする。
命題\ref{prop:rank-fifteen-pencil}の種数5曲線は、$[K:T:H:A:B]$ を座標とする
$\mathbb P^4$ の三つの二次式
\[
 H^2=K^2-T^2,\qquad A^2=(K+sT)^2-T^2,\qquad B^2=(K-sT)^2-T^2
\]
の滑らかな完全交叉である。二位置移動による商は $[H:A:B]$ を座標とする
滑らかな平面四次曲線で、その式は $D=A^2+B^2-2H^2$ と置いて
\[
 s^2(A^2-B^2)^2=4D(D+2s^2H^2)
\]
となる。
\end{proposition}
\begin{proof}
$[K:T]$ の各点上で三つの根を取るので射は有限であり、
三二次式は余次元3の完全交叉を与える。分岐値が相異なるので、
$T\ne0$ で零になる平方根は高々一つで、その平方の $K/T$ における零点は単純である。
$T=0$ では三根が全て非零である。従ってJacobi行列の階数は常に3で、
命題\ref{prop:rank-fifteen-pencil}の既約な関数体を持つ滑らかな完備曲線を得る。
随伴公式から埋込みの超平面束は標準束である。

$H,A,B$ は同時に零にならず、射影は二位置移動による三根の同時反転で不変である。
$D=2s^2T^2$、$A^2-B^2=4sKT$ を代入すると四次式を得る。
一般の像点で $T^2,KT$ が決まるため、残る二点は同時反転で移り合う。
射影の超平面束の次数は8であり、像の次数は4となる。従って表示した四次式が像の方程式である。
正規化は不動点のない商なので種数3で、これは平面四次の算術種数に等しい。
よって像は滑らかである。
\end{proof}

\begin{theorem}\label{thm:rank-fifteen-pencil-unramified}
一般の種数5水準の両表示で、$K=0$ と $K^2+sK-1=0$ 上の有限非分岐境界24点には、
通常部分の全ての窓が標準Koszul左右AS正則となるZariski近傍がある。
具体的に $c$ が $3/4$ に還元される標数零の通常窓で、
101において $K=0$ に還元されるもの、および197において $K^2+sK-1=0$ に
還元されるものは、各始点で四次中心の多項式環上の階数256の自由加群である。
\end{theorem}
\begin{proof}
$T=1$ とし、$L^2=1+c^2$、$R=K+(K^2-1)/s$、$W=K-(K^2-1)/s$ を取る。
$n=H/L$、$r=R\pm\sqrt{R^2-1}$、$w=W\pm\sqrt{W^2-1}$ で窓を表す。
$z=r^2$、$\lambda=w/r$ である。$R=0$ が上記の二次条件に当たる。
一般の $c$ では指定した三つの $K$ が相異なり、各々で三つの平方根は全て非零である。
関係の表示に必要な $z,\lambda,z-1,w^2-1$ は単元となる。
従ってこのチャートは境界そのものへ延びる。

$c=3/4$、$L=5/4$ と置く。$s=25/12$ で、101上の $K=0$ と
197上の二根 $K=35,45$ はいずれも上記の平方根を全て単純に分解する。
\path{computations/verify_rank_fifteen_pencil_unramified.py} は各根上の八通りすべてを取る。
各表示で計24点について隣接階数6、経路長4,5の次元35,56、
中心の文字次元 $(3,2,0,0)$、全四始点の四中心商の長さ13までの標準列を確認する。
$L$ の符号変更は $H$ の二択に含まれる。
命題\ref{prop:fixed-residual-regularity}の一般窓の経路・中心次元の下界と
剰余の上界が一致するので、中心を持ち上げ、四中心商の上界から
通常な持上げの全次数の自由性と正則性が従う。

各境界切断は $c$ の非空開集合の上で有限エタールである。
単純分解する上記の全点を検査したため、標数零でどの成分にも成功点がある。
道と中心の階数が一定の近傍で、四式の各始点の商の上界を与える小行列式の非消滅は開条件である。
有限な境界被覆から失敗集合の像を除けば、一般水準の24点を同時に含む所要の開近傍を得る。
\end{proof}
